\section{Introduction}

The deployment of Large Language Models (LLMs) is governed by a steep ``Intelligence Tax'': the assumption that higher performance requires linearly higher inference costs. Consequently, industrial applications often default to routing all traffic to frontier models (e.g., GPT-4-Turbo and its variants), treating cheaper models as strictly inferior approximations.

This assumption is increasingly flawed. Recent benchmarks reveal that specialized open-weights models (e.g., Mixtral-8x7B, Llama-3) can match or exceed the performance of frontier models on specific task clusters while costing orders of magnitude less. This phenomenon, which we term \emph{Quality Inversion}, transforms routing from a simple interpolation problem into a complex discovery problem. The goal is no longer just to ``save money on easy queries,'' but to identify the sparse, non-obvious sub-manifolds where the expensive model provides unique value.

Current state-of-the-art routers, such as RouteLLM~\cite{ong2024routellm} and FrugalGPT~\cite{chen2023frugalgpt}, address this via static supervision. They train a ``hardness classifier'' on offline datasets to predict which model to use. However, these static approaches suffer from two critical failures in production:

\begin{itemize}
    \item \textbf{Prior Rigidity}: They are trained on general-purpose public datasets (e.g., LMSYS Chatbot Arena~\cite{zheng2023judging}) that may not reflect domain-specific quality inversions. As we demonstrate, when a static router is forced to utilize an expensive model based on general ``hardness'' features, it often degrades performance on tasks where the cheaper model is actually superior.
    
    \item \textbf{The ``Cold Start'' Catastrophe}: The LLM landscape is highly non-stationary; new state-of-the-art models are released weekly. Static routers require expensive re-training and data collection to admit new models. Standard contextual bandits offer an online alternative but suffer from ``cold start'' periods where they must randomly explore (and fail) on thousands of requests before learning the utility of a new model.
\end{itemize}

In this work, we introduce \textbf{banditGPT}, a cost-aware routing framework that reformulates LLM selection as a \emph{Lifelong Learning} problem. Instead of a static classifier, we propose a hierarchical \emph{Expert Corralling} architecture~\cite{agarwal2017corralling} that manages a portfolio of routing strategies. To ensure safety, the router utilizes a mixing-enabled update rule that allows it to ``decisively decommission'' priors that are invalidated by online evidence. To ensure agility, we introduce \emph{Latent Semantic Transfer}, a protocol that initializes new models using semantic neighbors.

Through comprehensive experimental validation, we establish a \textbf{three-regime framework} that unifies adaptation behavior through learning rate selection: (1) \emph{Cold-start regime} ($\eta=0.1$--$0.3$): exploits semantic transfer for immediate deployment benefit via implicit regularization; (2) \emph{Pareto regime} ($\eta=1.0$): balances cost-quality trade-offs with partial adaptation; (3) \emph{Convergence regime} ($\eta=2.0$--$5.0$): complete prior unlearning and optimal policy convergence; (4) \emph{Safety regime} ($\eta=0.3$--$1.0$): catastrophic failure detection within 3--50 steps. Critically, we validate that semantic transfer operates through \textbf{implicit regularization} (26$\times$ more initial variance than cold start) rather than semantic accuracy---semantic similarity does not predict task-level performance correlation. This honest mechanism understanding explains both short-term benefit (14\% boost) and long-term need for aggressive unlearning, with timescale separation ensuring safety even when semantic priors are directionally wrong.

\subsection{Contributions}

Our contributions are as follows:

\begin{enumerate}
    \item \textbf{Empirical Analysis of Quality Inversion}: Adopting the standard RouteLLM evaluation topology (Mixtral-8x7B vs.\ GPT-4-Turbo) to ensure strict benchmarking consistency, we demonstrate a real-world scenario where the cheaper model outperforms the frontier baseline on average ($0.823$ vs.\ $0.812$). This isolates the algorithmic advantage of our router and proves that static ``intelligence tax'' assumptions are detrimental even when using optimized frontier models.
    
    \item \textbf{Three-Regime Framework}: We establish a unified design space for adaptive routing through comprehensive ablation across learning rates $\eta \in \{0.1, 0.3, 0.5, 1.0, 2.0, 5.0\}$. Rather than contradictory results, our experiments demonstrate complementary operating regimes: cold-start ($\eta=0.1$--$0.3$, short-term prior exploitation), Pareto sweep ($\eta=1.0$, balanced cost-quality), safety ($\eta=0.3$--$1.0$, catastrophic detection), and convergence ($\eta=2.0$--$5.0$, complete prior unlearning). This framework enables practitioners to select learning rates matching operational objectives (rapid deployment vs quality maximization vs safety-critical).
    
    \item \textbf{Algorithm (Expert Corralling)}: We implement a robust bandit architecture with a mixing parameter ($\gamma=0.05$) that prevents ``expert death,'' allowing the system to recover from non-stationary distribution shifts. Through rigorous multi-seed evaluation (N=10--20 seeds) with statistical significance testing (t-tests, Mann-Whitney U, Bonferroni correction), we validate catastrophic failure detection within 3--50 steps (Safety Regime, $\eta=0.3$) while enabling complete prior unlearning by step 1,121 (Convergence Regime, $\eta=5.0$). Timescale separation ($10\times$ faster detection than unlearning) ensures safety even when priors are directionally wrong.
    
    \item \textbf{Algorithm (Semantic Transfer Mechanism)}: We demonstrate that semantic transfer operates through \textbf{implicit regularization} (symmetry breaking via 26$\times$ more initial variance) rather than semantic accuracy. Direct measurement reveals semantic similarity does \emph{not} predict task-level performance correlation ($r=-0.38$, $p=0.75$). This honest mechanism understanding explains both short-term benefit (14\% boost via Conservative Learning Regime, $\eta=0.1$) and need for aggressive unlearning (Convergence Regime, $\eta=5.0$) for long-term optimality, transforming semantic transfer from a ``magic'' heuristic into a principled regularization technique.
    
    \item \textbf{Performance}: We achieve a state-of-the-art reward of $0.912 \pm 0.006$ on production traffic, closing 68.5\% of the optimality gap while reducing costs by 27\% compared to the GPT-4-Turbo baseline. We validate operational flexibility: tabula rasa (0.923) outperforms hybrid (0.912) with moderate learning ($\eta=1.0$), demonstrating the ``partial adaptation trap'' when priors are mismatched and adaptation time insufficient. Under severe domain mismatch with appropriate learning rate selection, Corralling achieves competitive performance (1.03--1.20$\times$ relative to baseline) while providing strong safety guarantees (39--43\% improvement over harmful warmup priors).
\end{enumerate}

\subsection{Precision in Baseline Specification}

Throughout this paper, we explicitly distinguish between the original GPT-4 model (released March 2023) and GPT-4-Turbo (the optimized variant with improved efficiency). Our empirical evaluation uses \textbf{GPT-4-Turbo} as the frontier baseline, reflecting the current production standard. When we refer to ``GPT-4 and its variants'' in general discussion, we are describing the broader class of frontier models; all quantitative comparisons are against the specific GPT-4-Turbo deployment configuration used in our dataset.

